% !tp-example Basic usage
\documentclass{article}
\begin{document}
\wcTpLatexViewer{
  \wcScript[type=tp/latex]|\section{Hello}A formula $x^2$.|
}
\end{document}

% !tp-example Javascript API
\documentclass{article}
\begin{document}
\wcScript[type=module]|
import "/components/latex/latex.js";
import "/components/latex-viewer/latex-viewer.js";

const example = document.createDocumentFragment();

const element1 = document.createElement("tp-latex-viewer");

const element2 = document.createElement("script");
element2.setAttribute("type", "tp/latex");
element2.textContent = "\\section{Hello}A formula $x^2$.";
element1.append(element2);
example.append(element1);

document.body.append(example);
|
\end{document}

% !tp-example Compare presentation modes
\wcSection{
  \subsection*{mode="default" — Standard controls}
  Open Code to edit the source, then click Run or press \wcKbd*{Shift}+\wcKbd*{Enter} in the editor.

  \wcTpLatexViewer[mode=default]{
    \wcScript[type=tp/latex]|
\section{Viewer modes}
Change this \textbf{message}, then run the source.
  |
  }
}

\wcSection{
  \subsection*{mode="lite" — Compact controls}
  Open Code to edit the source, then click Run or press \wcKbd*{Shift}+\wcKbd*{Enter} in the editor.

  \wcTpLatexViewer[mode=lite]{
    \wcScript[type=tp/latex]|
\section{Viewer modes}
Change this \textbf{message}, then run the source.
  |
  }
}

\wcSection{
  \subsection*{mode="simple" — Editor and rendered output only}
  Edit the code below, keep the cursor in the editor and press \wcKbd*{Shift}+\wcKbd*{Enter} (Return) to run it. This viewer has no Run button.

  \wcTpLatexViewer[mode=simple]{
    \wcScript[type=tp/latex]|
\section{Viewer modes}
Change this \textbf{message}, then run the source.
  |
  }
}
